\documentclass[11pt]{article}
\usepackage[a4paper,margin=0.85in]{geometry}
\usepackage[T1]{fontenc}
\usepackage{lmodern}
\usepackage{microtype}
\usepackage{graphicx}
\usepackage{booktabs}
\usepackage{tabularx}
\usepackage{array}
\usepackage{enumitem}
\usepackage{amsmath}
\usepackage{url}
\usepackage[hidelinks]{hyperref}
\usepackage[numbers,sort&compress]{natbib}
\usepackage{titlesec}
\setlength{\parindent}{1em}
\setlength{\parskip}{2pt}
\setlist{nosep,leftmargin=*}
\titleformat{\section}{\bfseries\large}{\thesection.}{0.5em}{}
\titleformat{\subsection}{\bfseries\normalsize}{\thesubsection}{0.5em}{}
\title{Cybersecurity in Connected Vehicles: Verification of Key Vulnerabilities, Assessment of Existing Tools, and a Layered Security Approach}
\author{Student Name\\
Department / Institution\\
Course: Cybersecurity for Connected Vehicles}
\date{}

\begin{document}
\maketitle

\begin{abstract}
Connected vehicles combine in-vehicle electronic control units (ECUs), wireless interfaces, cloud services, mobile applications, and vehicle-to-everything (V2X) communication. This connectivity improves safety, navigation, diagnostics, and user convenience, but it also expands the attack surface. A weakness in a remotely reachable component may expose vehicle data or, in severe cases, affect safety-critical functions. This article reviews major connected-vehicle security risks and examines how they can be verified using established defensive tools and controlled laboratory methods. The review covers Controller Area Network (CAN) message injection and replay, V2X message manipulation, telematics and infotainment exposure, Global Navigation Satellite System (GNSS) spoofing, and insecure over-the-air (OTA) updates. It compares representative market tools, including Wireshark, Zeek, Snort, Suricata, CANoe, and vulnerability scanners, while explaining why no single product provides complete end-to-end protection. A layered approach is proposed that combines asset inventory, threat modelling, authenticated communication, secure boot and signed updates, network segmentation, intrusion detection, privacy controls, and continuous incident response. A small, non-invasive Python example is included to demonstrate baseline CAN-log anomaly checks on an offline CSV file. The article argues that effective vehicle cybersecurity requires lifecycle governance and coordinated verification across the vehicle, communication infrastructure, and backend services rather than reliance on a single detection product.
\end{abstract}

\noindent\textbf{Keywords:} Connected vehicles, automotive cybersecurity, CAN bus, V2X, intrusion detection, OTA updates, ISO/SAE 21434, UN R155, vehicle security testing.

\section{Introduction}
Modern vehicles are increasingly software-defined cyber-physical systems. A typical connected vehicle may contain dozens of ECUs, sensors, infotainment equipment, cellular telematics, Bluetooth and Wi-Fi interfaces, GNSS receivers, and links to manufacturer cloud platforms. Advanced driver-assistance systems (ADAS), remote diagnostics, smartphone keys, fleet-management services, and V2X applications depend on the exchange of data across these components. The resulting connectivity creates benefits, but it also introduces interfaces that may be misconfigured, insufficiently authenticated, or vulnerable to software defects.

Unlike conventional information technology, a vehicle is a safety-critical environment. A cyber incident can affect confidentiality and privacy, but it may also affect the availability or integrity of functions such as braking assistance, steering support, propulsion management, or collision warnings. The impact depends on the architecture and the degree of isolation between the compromised component and safety-critical networks. It is therefore inaccurate to assume that every infotainment compromise automatically enables control of a vehicle; verification must establish the actual attack path and reachable functions.

The attack surface can be considered in three connected layers. The first is the in-vehicle layer, where ECUs communicate over networks such as CAN and Automotive Ethernet. The second is the external communication layer, including V2X, cellular, Wi-Fi, Bluetooth, GNSS, and diagnostic interfaces. The third is the backend ecosystem, including mobile applications, cloud APIs, update servers, identity systems, and fleet platforms. Security gaps at the boundaries between these layers are especially important because data may cross trust zones.

This paper has four objectives: (1) identify representative security problems in connected vehicles; (2) review research and standards relevant to those problems; (3) evaluate the role and limitations of existing market tools; and (4) propose a practical, layered verification and mitigation approach. The focus is defensive. Testing should be performed only on a vehicle or bench setup for which the tester has explicit authorization. No experiment described here requires transmitting attack traffic to a public road vehicle.

\section{Literature Review}
Automotive cybersecurity research has developed from early demonstrations of weaknesses in vehicle electronics to broader work on V2X trust, software updates, lifecycle risk management, and continuous monitoring. Checkoway et al. demonstrated that vulnerabilities in external interfaces and embedded systems could be chained to reach vehicle networks, highlighting the importance of considering the complete attack path rather than evaluating components in isolation \cite{checkoway2011}. Miller and Valasek later reported remote attack paths involving vehicle connectivity and internal networks, reinforcing the need for segmentation and secure gateway design \cite{miller2015}.

Recent surveys address the increased complexity of connected and automated vehicles. Yoshizawa et al. review security and privacy challenges in V2X systems and identify weaknesses that can arise across communication layers, standards, and deployment models \cite{yoshizawa2023}. Their work highlights that cryptographic protection alone does not eliminate threats such as compromised legitimate participants, privacy leakage, or inconsistent implementation. A 2024 review by Alshammari et al. surveys automotive cybersecurity frameworks, standards, testing, and monitoring technologies, and discusses the challenges of integrating security into complex vehicle software ecosystems \cite{alshammari2024}. A separate 2024 analysis of connected-vehicle cybersecurity examines attacks on V2X traffic and trust-based approaches, including replay, denial-of-service, and false-message injection \cite{cyacv2024}.

Software update security is another major research area. Li et al. survey OTA upgrade techniques for intelligent connected vehicles and describe risks across cloud infrastructure, update terminals, and vehicle-side update objects \cite{li2024ota}. An update pipeline can become a high-impact entry point if packages are not authenticated, signing keys are poorly protected, rollback is possible, or update delivery infrastructure is compromised. Secure update design therefore needs more than encrypted transport: the vehicle must verify the authenticity and integrity of the package, enforce version and compatibility policy, and recover safely from failed installation.

Standards and regulatory guidance translate these research concerns into engineering and organizational processes. ISO/SAE 21434 defines cybersecurity engineering activities across the road-vehicle lifecycle, including risk assessment and treatment \cite{iso21434}. UN Regulation No. 155 establishes requirements for a Cyber Security Management System and vehicle-type cybersecurity, while UN Regulation No. 156 addresses software update management systems \cite{unece155,unece156}. NHTSA's 2022 best-practice guidance recommends risk-based protection, vulnerability identification, testing, monitoring, containment, and remediation \cite{nhtsa2022}. These sources collectively support a lifecycle approach: security must be considered during design, production, operation, maintenance, and decommissioning.

\subsection{Research gap}
The literature and standards provide extensive guidance, but practical deployments still face integration challenges. Tools often specialize in one layer: packet analyzers inspect network traffic, CAN tools decode vehicle bus data, vulnerability scanners examine hosts and services, and endpoint monitoring detects selected anomalies. These tools may not share a common vehicle asset model, threat context, or response workflow. Furthermore, a detection alert is not equivalent to a verified vulnerability. A disciplined process must reproduce a finding in a safe test environment, determine its preconditions and impact, document evidence, and verify that the mitigation actually closes the issue without affecting vehicle safety or performance.

\section{Major Security Problems in Connected Vehicles}
\subsection{In-vehicle CAN communication}
CAN was designed for reliable real-time communication among ECUs, not as a modern security protocol with built-in sender authentication or confidentiality. In many deployments, a node able to access the bus can transmit frames using an identifier that other nodes may accept. Risks include message spoofing, replay, bus flooding, and unauthorized diagnostic requests. The practical impact depends on gateway filtering, network topology, message plausibility checks, and the functions controlled by receiving ECUs. Defences include secure gateways, network segmentation, authenticated diagnostic access, message freshness mechanisms where feasible, and intrusion detection based on timing and signal behaviour.

\subsection{V2X communication threats}
V2X enables vehicle-to-vehicle, vehicle-to-infrastructure, and vehicle-to-network communication. It supports cooperative awareness, hazard notifications, and traffic coordination. Attackers may attempt replay, false-data injection, identity misuse, denial-of-service, or tracking through persistent identifiers. Digital certificates and signed messages can establish message origin and integrity, but they do not guarantee that a credentialed sender is honest or that sensor data is accurate. Certificate lifecycle management, misbehaviour detection, privacy-preserving pseudonyms, and plausibility checks are therefore also required \cite{yoshizawa2023}.

\subsection{Telematics, infotainment, and external interfaces}
Cellular modems, companion applications, Bluetooth, Wi-Fi, USB, and diagnostic ports provide useful functionality but create potential entry points. Weak authentication, exposed services, insecure API authorization, outdated libraries, and poor separation between infotainment and safety-related networks can increase risk. Verification should inventory interfaces, inspect service exposure, review application permissions and API access controls, and confirm that gateways restrict traffic to explicitly required destinations and message types.

\subsection{GNSS spoofing and sensor manipulation}
Connected vehicles use GNSS for navigation, fleet tracking, timing, and sometimes as an input to automated functions. Spoofed or jammed satellite signals can produce misleading position or timing information. Sensor manipulation can also target cameras, radar, lidar, or data-fusion inputs. Mitigation should avoid trusting a single sensor: compare GNSS with inertial measurements, wheel odometry, map constraints, and other independent sources; monitor signal-quality indicators; and transition to a safe operating mode when confidence is insufficient.

\subsection{Insecure OTA updates and supply-chain weaknesses}
OTA updates are essential for patching vulnerabilities, but the update process itself must be protected. Threats include compromised update servers, stolen signing keys, tampered packages, downgrade attacks, interrupted updates, and vulnerable third-party components. A robust design uses signed packages, protected key storage, mutual authentication, secure boot, anti-rollback controls, staged deployment, software bills of materials (SBOMs), and a recovery mechanism. The vehicle should verify updates locally rather than trusting that a package is safe merely because it arrived over an encrypted connection \cite{li2024ota}.

\subsection{Privacy and data protection}
Vehicles may collect location history, driving behaviour, contacts, voice data, diagnostic records, and usage patterns. Excessive collection, long retention periods, weak access controls, or insecure cloud APIs can expose sensitive information. Privacy controls should apply data minimization, purpose limitation, encryption in transit and at rest, access logging, retention limits, and clear user consent mechanisms.

\section{Verification Methodology and Existing Market Tools}
A safe verification process begins with written authorization, a defined scope, and a test environment. A bench setup with a CAN interface, simulator, or hardware-in-the-loop platform is preferable to experimentation on a road vehicle. The following workflow is suitable for a course-level assessment.

\begin{enumerate}
\item \textbf{Asset and interface inventory:} document ECUs, software versions, network segments, wireless interfaces, cloud endpoints, diagnostic services, and data flows.
\item \textbf{Threat modelling:} identify assets, trust boundaries, attacker capabilities, attack paths, and safety/privacy consequences. Prioritize findings by likelihood and impact.
\item \textbf{Configuration and exposure review:} inspect service exposure, authentication policy, certificates, update configuration, API authorization, and software-component versions.
\item \textbf{Controlled traffic analysis:} capture authorized test traffic and inspect protocol behaviour, unexpected destinations, unusual message rates, and malformed frames.
\item \textbf{Vulnerability validation:} reproduce suspected weaknesses on a simulator or isolated bench, record evidence, and avoid unsafe actuation.
\item \textbf{Mitigation and retest:} apply a control, repeat the same test, check for regressions, and document residual risk.
\end{enumerate}

\begin{table*}[t]
\centering
\caption{Representative tools and their roles in connected-vehicle security}
\small
\begin{tabularx}{\textwidth}{p{2.3cm}p{3.2cm}Xp{3.0cm}}
\toprule
\textbf{Tool} & \textbf{Primary use} & \textbf{Strength in vehicle security} & \textbf{Important limitation}\\
\midrule
Wireshark & Packet capture and protocol analysis & Helps inspect Ethernet, IP, TLS metadata, and supported V2X-related traffic dissectors & Does not automatically understand all proprietary vehicle signals or prove exploitability\\
Zeek & Network security monitoring & Produces structured logs and supports behavioural analysis of IP-based vehicle/backend traffic & Requires suitable traffic visibility and custom policies for automotive protocols\\
Snort / Suricata & Network intrusion detection and prevention & Signature and rule-based alerts for suspicious network activity & Encrypted traffic and unknown attacks can reduce visibility; tuning is essential\\
Vector CANoe / CANalyzer & Automotive network simulation, testing, and analysis & Supports CAN, LIN, Ethernet, diagnostics, and ECU test workflows depending on configuration & Commercial licensing and specialist configuration may be required\\
Nessus / OpenVAS & Host and service vulnerability assessment & Identifies known vulnerabilities and insecure configurations in reachable infrastructure & Not a substitute for ECU-specific analysis or vehicle safety validation\\
Burp Suite & Web and API security testing & Useful for authorized testing of companion apps and cloud-facing web APIs & Does not assess in-vehicle buses or radio-layer V2X security\\
\bottomrule
\end{tabularx}
\end{table*}

The tools are complementary rather than interchangeable. Wireshark is useful for manual packet inspection; Zeek provides network telemetry; Snort and Suricata detect traffic patterns matching configured rules; CANoe and CANalyzer support automotive bus simulation and testing; vulnerability scanners identify known weaknesses in reachable hosts; and Burp Suite helps assess web applications and APIs. A tool's presence does not prove that a vehicle is secure. Coverage depends on sensor placement, protocol support, configuration quality, access to relevant logs, and the ability to validate alerts.

\section{Technology Used to Solve the Problems}
\subsection{Secure architecture and segmentation}
A secure vehicle architecture separates infotainment, telematics, body, powertrain, and safety-critical domains using gateways and policy enforcement. Gateways should apply allowlists based on message identifiers, source, destination, rate, and operating context. Diagnostic access should require strong authorization and be limited by role and vehicle state. Segmentation reduces the chance that compromise of a less-trusted component becomes access to safety-critical functions.

\subsection{Cryptography and identity}
Use modern, well-reviewed cryptographic protocols rather than proprietary encryption. Mutual authentication can protect communication between vehicles, roadside infrastructure, mobile applications, and backend services. V2X systems can use certificate-based signatures and pseudonym certificates to provide message integrity while reducing long-term linkability. Keys should be provisioned, stored, rotated, revoked, and retired through a managed lifecycle. Hardware security modules or trusted execution environments can protect sensitive keys where supported.

\subsection{Secure boot and authenticated OTA}
Secure boot verifies each software stage before execution. OTA packages should be digitally signed, checked for integrity, bound to the intended vehicle/ECU and version, and protected against rollback. Update servers should use strong access control, key separation, audit logs, and staged release processes. A/B partitions or another validated recovery design can reduce the risk of an interrupted update leaving an ECU unusable. These controls align with the lifecycle focus of UN R155 and R156 \cite{unece155,unece156}.

\subsection{Intrusion detection and anomaly monitoring}
An in-vehicle intrusion detection system (IDS) can monitor CAN identifiers, message frequency, payload changes, diagnostic activity, and cross-network flows. Rule-based detection is interpretable and efficient for known violations, while statistical or machine-learning methods may identify deviations from learned normal behaviour. However, a model can generate false positives when driving conditions change, and an attacker may attempt to mimic normal traffic. A practical design combines deterministic safety rules, contextual baselines, model monitoring, and human review. Detection should trigger a proportionate response, such as logging, alerting, isolating a non-critical interface, or requesting a safe-state transition under validated vehicle policy.

\subsection{Privacy-preserving data governance}
Collect only data needed for a defined purpose, apply retention limits, and separate identifiers from operational telemetry where possible. Encrypt data in transit and at rest, enforce least-privilege access, and maintain auditable access logs. For fleet analytics, aggregation or privacy-preserving learning can reduce the need to centralize raw driver-level data. Privacy requirements must be considered alongside safety and operational needs.

\section{Proposed Layered Solution and Comparative Analysis}
This article proposes a layered vehicle-security verification framework. It is not a claim that one new product can replace the market. Instead, it combines established tools with controls that address their gaps.

\begin{table}[t]
\centering
\caption{Security problem, verification method, and recommended mitigation}
\small
\begin{tabularx}{\columnwidth}{p{1.7cm}p{2.1cm}X}
\toprule
\textbf{Problem} & \textbf{Verification} & \textbf{Mitigation}\\
\midrule
CAN spoofing/replay & Offline log review and isolated bench tests & Gateway filtering, freshness checks, IDS\\
V2X false messages & Certificate and plausibility checks in a test network & PKI lifecycle, misbehaviour detection\\
Exposed services & Authorized scanning and configuration review & Reduce attack surface, harden services\\
OTA tampering & Signature, version, and rollback tests & Signed packages, secure boot, anti-rollback\\
GNSS spoofing & Simulator-based signal and sensor consistency tests & Multi-sensor fusion and confidence monitoring\\
Data leakage & API authorization and privacy review & Minimization, encryption, least privilege\\
\bottomrule
\end{tabularx}
\end{table}

Compared with using a single market tool, the proposed approach has four advantages. First, it maps each tool to the layer it can actually observe, reducing false assumptions about coverage. Second, it combines technical testing with threat modelling and lifecycle governance. Third, it treats mitigation verification as a required step rather than stopping at an alert. Fourth, it includes privacy and software-supply-chain concerns alongside network attacks.

A practical improvement over isolated tools would be a unified evidence workflow. Each finding should include the affected asset and software version, the test condition, timestamped evidence, severity rationale, owner, mitigation, and retest result. An SBOM can help map a newly disclosed component vulnerability to affected ECUs and vehicle versions. A centralized security operations process can correlate vehicle telemetry, cloud logs, certificate events, and update status. Such integration improves prioritization, but it must preserve data minimization and avoid creating an overly privileged monitoring system.

\section{Illustrative Offline CAN-Log Verification}
To make the methodology reproducible, the accompanying repository contains a small Python script, \texttt{src/can\_log\_checker.py}, that analyses an offline CSV file. It checks whether CAN identifiers belong to a configured allowlist, whether the data length code (DLC) is outside the conventional range of 0--8 for classical CAN, and whether an identifier exceeds a configurable message-rate threshold. The script does not connect to a vehicle, transmit frames, decode proprietary signals, or prove that a detected event is malicious. It is a simple baseline suitable for demonstrating the verification workflow.

The sample dataset is synthetic and intentionally includes a small number of anomalous rows. In a real assessment, the allowlist and rate thresholds must be derived from the specific vehicle architecture and validated operating conditions. Classical CAN and CAN FD have different payload limits; the example deliberately uses the classical CAN limit and should be adjusted for CAN FD logs.

\section{Discussion}
No existing market tool can independently guarantee connected-vehicle security. Network tools may miss in-vehicle signal semantics; automotive bus tools may not assess cloud APIs; vulnerability scanners may not understand safety implications; and IDS products may generate alerts without enough context to determine whether an event is hazardous. Proprietary protocols, encrypted links, incomplete asset inventories, and limited access to OEM data further complicate assessment.

The most important improvement is therefore not simply replacing one commercial product with another. It is to establish a repeatable engineering process in which tools are selected according to the asset and threat, tests are conducted safely, findings are validated, and controls are retested. The approach should be aligned with ISO/SAE 21434, UN R155/R156, and applicable national requirements. AI-based detection and blockchain-based trust mechanisms may be useful in selected scenarios, but they should not be adopted without measuring latency, computational overhead, false-positive rates, privacy impact, and failure behaviour. In a vehicle, a security mechanism that interferes with deterministic safety functions can create a different hazard.

\section{Conclusion}
Connected vehicles face security risks across in-vehicle networks, V2X communication, telematics, GNSS, OTA updates, cloud services, and personal data processing. Research and regulatory guidance show that these risks must be managed throughout the vehicle lifecycle. Existing tools such as Wireshark, Zeek, Snort, Suricata, CANoe, vulnerability scanners, and Burp Suite provide valuable but specialized capabilities. Their limitations are best addressed through coordinated use, a complete asset inventory, threat modelling, secure architecture, cryptographic identity, signed updates, intrusion detection, privacy governance, and incident response.

The proposed framework emphasizes controlled verification and retesting rather than relying on product claims or isolated alerts. The included offline CAN-log checker illustrates how basic evidence can be collected without interacting with a live vehicle. Future work could extend the demonstration with CAN FD support, signal-aware specifications, V2X certificate validation, and carefully evaluated machine-learning models using publicly available datasets. Any extension should report its dataset, test conditions, detection metrics, and limitations transparently. Ultimately, connected-vehicle cybersecurity is a continuous safety and engineering responsibility shared by manufacturers, suppliers, service providers, and vehicle operators.

\bibliographystyle{unsrtnat}
\bibliography{references}
\end{document}
